\chapter{Software Architecture}
\label{chap:architecture}

\section{Architecture Overview}
\label{sec:architecture-overview}\index{architecture}\index{login}\index{guest}\index{Keycloak}\index{Caddy}

A question crosses eight layers (Figure~\ref{fig:architecture}): from the web chat to access
control, the input guardrail, the LangGraph\index{LangGraph} orchestrator~\cite{langgraph}, the
sources, the generation and the output guardrail, up to the answer or an operator. Login is optional: a guest gets general answers from the
public documents, a customer also gets answers on their own policy. Everything is free: one
server, plus Groq's free \acs{api} for the two large models (Figure~\ref{fig:overview-strip}).

% Figure 2.1: a question at a glance, from the chat to the answer (details in the next figure)
\begin{figure}[H]
  \centering
  \begin{tikzpicture}[font=\scriptsize,
      lay/.style={rounded corners=3pt, text width=1.5cm, minimum height=1.55cm, align=center,
                  inner sep=2pt},
      flow/.style={-{Stealth[length=1.8mm]}, accent, line width=0.7pt}]
    \newcommand\lgx[1]{\includegraphics[height=0.5cm, width=0.62cm, keepaspectratio]{logos/#1}}
    % \lay{index}{category}{icon or logo}{label}
    \newcommand\lay[4]{\node[lay, cat #2] (l#1) at ({(#1-1)*1.98},0)
      {#3\\[3pt]{\bfseries #4}};}
    \lay{1}{infra}{{\color{accent}\Large\faComments}}{Web chat}
    \lay{2}{sec}{\lgx{keycloak}}{Optional login}
    \lay{3}{sec}{\lgx{qwen}}{Input guardrail}
    \lay{4}{infra}{\lgx{langgraph}}{Orche\-strator}
    \lay{5}{data}{\lgx{opensearch}\,\lgx{mcp}}{Sources}
    \lay{6}{model}{\lgx{openai}}{Generation}
    \lay{7}{sec}{{\color{catsec}\Large\faCheckDouble}}{Output guardrail}
    \lay{8}{infra}{\lgx{chatwoot}}{Answer or operator}
    \foreach \a/\b in {1/2,2/3,3/4,4/5,5/6,6/7,7/8}
      \draw[flow] (l\a) -- (l\b);
  \end{tikzpicture}
  \caption{A question at a glance. Figure~\ref{fig:architecture} shows each layer in detail.}
  \label{fig:overview-strip}
\end{figure}

% Figure 2.1: the eight layers crossed by a question, with the logo of every technology.
% Tools sit in three columns, at x = 4.85, 9.25 and 13.65.
\begin{figure}[p]
  \centering
  \begin{tikzpicture}[
      font=\scriptsize,
      layer/.style={draw=black!30, fill=black!4, rounded corners=4pt},
      band/.style={draw=black!45, dashed, rounded corners=4pt},
      tool/.style={rounded corners=2pt, minimum height=0.85cm, inner sep=2pt},
      layername/.style={anchor=west, text width=2.3cm, align=left, font=\footnotesize\bfseries},
      grouptitle/.style={font=\scriptsize\bfseries, text=black!65},
      hop/.style={-{Stealth[length=2mm]}, black!55, line width=0.6pt},
      hoplabel/.style={anchor=west, font=\scriptsize\itshape, text=black!70},
      swatch/.style={rounded corners=1.5pt, minimum width=0.55cm, minimum height=0.32cm},
    ]
    % \tool{category}{x}{y}{logo or icon}{name}{role}
    \newcommand\tool[6]{\node[tool, cat #1] at (#2,#3) {%
      \begin{tabular}{@{}>{\centering\arraybackslash}m{0.55cm}@{\hspace{3pt}}
                         >{\raggedright\arraybackslash}m{3.45cm}@{}}
        #4 & \textbf{#5}\newline #6
      \end{tabular}};}
    \newcommand\lgo[1]{\logo[0.5cm]{#1}}
    \newcommand\ico[1]{{\color{accent}\Large #1}}
    % \layer{top}{bottom}{number and name}
    \newcommand\layer[3]{%
      \draw[layer] (0,#1) rectangle (15.85,#2);
      \node[layername] at (0.15,{(#1+#2)/2}) {#3};}
    % \hop{from y}{to y}{what travels between two layers}
    \newcommand\hop[3]{%
      \draw[hop] (0.9,#1) -- (0.9,#2);
      \node[hoplabel] at (1.05,{(#1+#2)/2}) {#3};}

    \layer{0.6}{-0.6}{1\enspace Channels}
    \tool{infra}{4.85}{0}{\ico{\faComments}}{Web chat}{HTML + JS, served by FastAPI}
    \hop{-0.6}{-1.1}{message, + JWT after login}

    \layer{-1.1}{-2.3}{2\enspace Access}
    \tool{infra}{4.85}{-1.7}{\lgo{caddy}}{Caddy}{HTTPS, reverse proxy}
    \tool{sec}{9.25}{-1.7}{\lgo{keycloak}}{Keycloak}{optional login, token exchange}
    \hop{-2.3}{-2.8}{verified identity}

    \layer{-2.8}{-4.0}{3\enspace Input guardrail}
    \tool{model}{4.85}{-3.4}{\lgo{qwen}}{Qwen3Guard-Gen-0.6B}{jailbreaks, banned content}
    \tool{sec}{9.25}{-3.4}{\lgo{presidio}}{Presidio}{masks personal data}
    \hop{-4.0}{-4.5}{admitted question}

    \layer{-4.5}{-5.7}{4\enspace Orchestrator}
    \tool{infra}{4.85}{-5.1}{\lgo{langgraph}}{LangGraph}{graph of 13 nodes}
    \tool{infra}{9.25}{-5.1}{\lgo{fastapi}}{FastAPI + SSE}{API, streamed states}
    \tool{model}{13.65}{-5.1}{\lgo{openai}}{gpt-oss-20b}{intent, rewording \textbullet\ Groq}
    \hop{-5.7}{-6.2}{query, filters, service token}

    \layer{-6.2}{-9.6}{5\enspace Sources\\[2pt]{\mdseries\color{black!60}read in parallel}}
    \node[grouptitle] at (4.85,-6.45) {Knowledge};
    \node[grouptitle] at (9.25,-6.45) {Customer data};
    \node[grouptitle] at (13.65,-6.45) {Memory};
    \tool{data}{4.85}{-7.125}{\lgo{opensearch}}{OpenSearch}{BM25 + vectors, filters}
    \tool{model}{4.85}{-8.075}{\lgo{qwen}}{Qwen3-Embedding-0.6B}{clause vectors}
    \tool{model}{4.85}{-9.025}{\lgo{qwen}}{Qwen3-Reranker-0.6B}{reorders the top 40}
    \tool{data}{9.25}{-7.125}{\lgo{mcp}}{MCP policy server}{after login, read only}
    \tool{infra}{9.25}{-8.075}{\ico{\faBuilding}}{Core insurance system}{existing company APIs}
    \tool{data}{13.65}{-7.125}{\lgo{postgresql}}{PostgreSQL}{checkpoints, audit}
    \tool{data}{13.65}{-8.075}{\lgo{valkey}}{Valkey}{cache, Chatwoot queues}
    \hop{-9.6}{-10.1}{clauses + policy data}

    \layer{-10.1}{-11.3}{6\enspace Generation}
    \tool{model}{4.85}{-10.7}{\lgo{openai}}{gpt-oss-120b}{117B MoE \textbullet\ Groq}
    \tool{model}{9.25}{-10.7}{\lgo{ollama}}{Qwen3-4B on Ollama}{last fallback, CPU}
    \hop{-11.3}{-11.8}{draft with citations}

    \layer{-11.8}{-13.0}{7\enspace Output guardrail}
    \tool{sec}{4.85}{-12.4}{\ico{\faCheckDouble}}{Validator}{citations, amounts, dates}
    \tool{model}{9.25}{-12.4}{\lgo{openai}}{gpt-oss-20b}{groundedness, criteria}
    \hop{-13.0}{-13.5}{verified answer}

    \layer{-13.5}{-14.7}{8\enspace Output}
    \tool{infra}{4.85}{-14.1}{\ico{\faIcon{file-alt}}}{Answer with sources}{citations, intent, confidence}
    \tool{infra}{9.25}{-14.1}{\lgo{chatwoot}}{Chatwoot}{operator console, corrections}

    % Services used by every layer
    \draw[band] (0,-15.3) rectangle (15.85,-18.7);
    \node[layername] at (0.15,-17.0)
      {Cross-cutting services\\[2pt]{\mdseries\color{black!60}used by every layer}};
    \node[grouptitle] at (4.85,-15.55) {Observability};
    \node[grouptitle] at (9.25,-15.55) {Quality};
    \node[grouptitle] at (13.65,-15.55) {Hosting};
    \tool{infra}{4.85}{-16.225}{\ico{\faStream}}{JSON logs}{a trace per node, masked}
    \tool{sec}{4.85}{-17.175}{\ico{\faLink}}{Audit chain}{hash-linked rows}
    \tool{infra}{9.25}{-16.225}{\logo[0.32cm]{ragas}\,\logo[0.32cm]{gemini}}{Ragas + Gemini}
      {judge of the tests}
    \tool{sec}{9.25}{-17.175}{\lgo{promptfoo}}{promptfoo}{red teaming, injection}
    \tool{infra}{9.25}{-18.125}{\lgo{pytest}}{pytest}{routes, scenarios}
    \tool{infra}{13.65}{-16.225}{\lgo{docker}}{Docker Compose}{every service, one file}
    \tool{infra}{13.65}{-17.175}{\lgo{oracle}}{Oracle Cloud}{free server, 4 cores, 24 GB}

    % Legend
    \node[swatch, cat model] at (0.4,-19.35) {};
    \node[anchor=west] at (0.75,-19.35) {open-weight model};
    \node[swatch, cat sec] at (3.8,-19.35) {};
    \node[anchor=west] at (4.15,-19.35) {security and controls};
    \node[swatch, cat data] at (7.7,-19.35) {};
    \node[anchor=west] at (8.05,-19.35) {data and knowledge};
    \node[swatch, cat infra] at (11.4,-19.35) {};
    \node[anchor=west] at (11.75,-19.35) {infrastructure and channels};
  \end{tikzpicture}
  \caption{The layers of the assistant, crossed by a question from top to bottom. Everything runs
    on one free server, except the two gpt-oss models (Groq's free \acs{api}).}
  \label{fig:architecture}
\end{figure}


\section{Components}
\label{sec:components}\index{policy server}\index{token exchange}\index{core insurance system}

Figure~\ref{fig:architecture} shows every component and its role; technologies, licences and
details are in Table~\ref{tab:components} (Appendix). The policy
server\index{Model Context Protocol}~\cite{mcp} exposes four read-only tools over the company's core
insurance system (Figure~\ref{fig:mcp-tools}). Keycloak issues its tokens
through standard token exchange~\cite{keycloak2026tokenexchange}.

% Figure 2.3: the four read-only tools of the policy server, one card each
\begin{figure}[H]
  \centering
  % \mcptool{icon}{function}{what it returns}{protection}
  \newcommand\mcptool[4]{%
    \begin{tcolorbox}[card]
      \cardtitle{#1}{\texttt{#2}}
      #3\par
      {\color{catsec}\faLock}\ #4
    \end{tcolorbox}}
  \begin{tcbraster}[raster columns=2, cardgrid, raster before skip=0pt]
    \mcptool{\faList}{list\_my\_policies()}{Product, \acs{cga} edition, status, dates, insured
      object masked.}{Only the token subject's policies.}
    \mcptool{\faFileContract}{get\_policy(policy\_id)}{Guarantees, limits, deductibles,
      premium.}{Someone else's policy: 404, not 403.}
    \mcptool{\faEuroSign}{get\_payments(policy\_id)}{Instalments, due dates, status.}
      {IBAN cut to its last four digits.}
    \mcptool{\faCarCrash}{get\_claims(policy\_id)}{Open claims, stage, contact.}
      {No reserves, no internal notes.}
  \end{tcbraster}
  \caption{The four read-only tools of the policy server, in front of the core insurance system.
    They are used only after login.}
  \label{fig:mcp-tools}
\end{figure}


\section{Component Interactions}
\label{sec:interactions}

The intent recognised by \code{classify} decides the route (Figure~\ref{fig:intents}). A turn is
one run of the graph (Figure~\ref{fig:graph}): after each node, a function written in code picks
the next one.\index{intent}

% Figure 2.3: the eight intents recognised by classify and the route each one takes
\begin{figure}[H]
  \centering
  \begin{tikzpicture}[
      font=\scriptsize,
      chip/.style={rounded corners=6pt, draw=catmodel!60, fill=catmodel!8, inner sep=3pt,
                   align=left, text width=2.75cm, anchor=west},
      dest/.style={rounded corners=3pt, draw=black!40, fill=white, inner sep=4pt, align=left,
                   text width=4.9cm},
      route/.style={-{Stealth[length=2mm]}, black!55, line width=0.6pt},
    ]
    \node[rounded corners=3pt, cat model, align=center, text width=1.7cm, minimum height=1.2cm]
      (cl) at (0.95,-2.5) {\logo[0.5cm]{openai}\\[2pt]{\ttfamily classify}};
    % \route{id}{y}{label}{meaning}{second label}{second meaning}{destination}
    \newcommand\intent[7]{%
      \node[chip] (a#1) at (2.5,#2) {{\ttfamily\bfseries #3}\\{\color{black!65}#4}};
      \ifx\relax#5\relax
        \coordinate (b#1) at (a#1.east);
      \else
        \node[chip] (c#1) at (5.75,#2) {{\ttfamily\bfseries #5}\\{\color{black!65}#6}};
        \coordinate (b#1) at (c#1.east);
      \fi
      \node[dest] (d#1) at (13.25,#2) {#7};
      \draw[route] (cl.east) -- (a#1.west);
      \draw[route] (b#1) -- (d#1.west);}
    \intent{1}{0}{polizza}{conditions of one's policy}{}{}
      {after login: \code{resolve_policy}, then \code{fetch_policy_data} and \code{retrieve} in
       parallel. A guest is asked to log in}
    \intent{2}{-1.3}{sinistro}{status of a claim}{pagamenti}{instalments, payments}
      {after login: \code{resolve_policy}, then \code{fetch_policy_data} with the claims or the
       instalments. A guest is asked to log in}
    \intent{3}{-2.6}{prodotto}{a product in general}{}{}
      {\code{retrieve} on the public documents, for guests and customers alike}
    \intent{4}{-3.8}{operatore}{asks for a person}{reclamo}{a complaint}
      {\code{handoff}; a complaint goes with high priority}
    \intent{5}{-5.0}{fuori\_ambito}{off-topic request}{saluto}{greeting}
      {\code{direct_reply}, a fixed text}
  \end{tikzpicture}
  \caption{The intents recognised by \texttt{classify} (labels as in the code) and their routes.}
  \label{fig:intents}
\end{figure}

% Figure 2.2: the LangGraph graph of a conversation turn; badges show the technology of each node
\begin{figure}[p]
  \centering
  \begin{tikzpicture}[
      font=\scriptsize,
      gnode/.style={rounded corners=3pt, text width=2.6cm, minimum height=0.8cm, align=center,
                    inner sep=2pt},
      term/.style={dashed},
      pill/.style={rounded corners=7pt, draw=black!55, fill=black!5, minimum width=1.4cm,
                   minimum height=0.45cm, font=\scriptsize\bfseries},
      route/.style={-{Stealth[length=2mm]}, black!60, line width=0.6pt},
      loop/.style={route, dashed},
      elabel/.style={font=\scriptsize, text=black!75, align=center, inner sep=1.5pt},
      badge/.style={circle, fill=white, draw=black!25, inner sep=1pt, minimum size=0.5cm},
    ]
    % \gnode{name}{category, extra styles}{x}{y}{node}{what it uses}
    \newcommand\gnode[6]{%
      \node[gnode, cat #2] (#1) at (#3,#4) {{\ttfamily #5}\\{\color{black!70}#6}};}
    % \badge{node}{logo}: the logo on the top-left corner of a node
    \newcommand\badge[2]{%
      \node[badge] at ([xshift=-1.6mm, yshift=1.6mm]#1.north west) {\includegraphics[height=0.3cm, width=0.34cm,
        keepaspectratio]{logos/#2}};}

    \node[pill] (start) at (4.6,0) {START};
    \gnode{ig}{sec}{4.6}{-1.2}{input\_guard}{Qwen3Guard 0.6B}
    \gnode{cl}{model}{4.6}{-2.75}{classify}{gpt-oss-20b}
    \gnode{dr}{infra, term}{0}{-2.75}{direct\_reply}{fixed replies}
    \gnode{rp}{data}{4.6}{-4.3}{resolve\_policy}{MCP, policy list}
    \gnode{ac}{infra, term}{0}{-4.3}{ask\_clarification}{which policy?}
    \gnode{fp}{data}{4.6}{-5.85}{fetch\_policy\_data}{MCP, read only}
    \gnode{rt}{data}{9.0}{-5.85}{retrieve}{search + rerank}
    \gnode{gr}{data}{4.6}{-7.4}{grade}{rerank threshold}
    \gnode{rq}{model}{9.0}{-7.4}{rewrite\_query}{gpt-oss-20b, glossary}
    \gnode{ge}{model}{4.6}{-8.95}{generate}{gpt-oss-120b}
    \gnode{ho}{infra, term}{12.9}{-8.95}{handoff}{Chatwoot + pause}
    \gnode{og}{sec}{4.6}{-10.5}{output\_guard}{validator + 20b}
    \gnode{fi}{infra}{4.6}{-12.05}{finalize}{sources, audit}
    \node[pill] (end) at (4.6,-13.3) {END};
    \badge{ig}{qwen} \badge{cl}{openai} \badge{rq}{openai}
    \badge{rp}{mcp} \badge{fp}{mcp} \badge{rt}{opensearch}
    \badge{ge}{openai} \badge{og}{openai} \badge{ho}{chatwoot}
    \badge{fi}{postgresql}

    \draw[route] (start) -- (ig);
    \draw[route] (ig) -- node[elabel, right] {admitted} (cl);
    \draw[route] (ig.west) -| node[elabel, above, pos=0.25] {blocked} (dr.north);
    \draw[route] (cl.west) -- node[elabel, above] {greeting,\\off-topic,\\no login} (dr.east);
    \draw[route] ([yshift=2mm]cl.east) -| (ho.north);
    \node[elabel, above] at (11.0,-2.55) {operator, complaint};
    \draw[route] ([yshift=-2mm]cl.east) -| (rt.north);
    \node[elabel, below] at (7.6,-2.95) {product};
    \draw[route] (cl) -- node[elabel, right] {policy, claim,\\payment} (rp);
    \draw[route] (rp.west) -- node[elabel, above] {no single\\policy} (ac.east);
    \draw[route] (rp) -- node[elabel, right] {policy found} (fp);
    \draw[route] (rp.east) -- node[elabel, above] {in parallel} (9.0,-4.3);
    \fill[black!60] (9.0,-4.3) circle (1.6pt);
    \draw[route] (fp) -- (gr);
    \draw[route] (fp.west) -- ++(-0.6,0) |- ([yshift=1.5mm]ge.west);
    \node[elabel, left] at (2.6,-7.4) {claim,\\payment};
    \draw[route] ([xshift=-6mm]rt.south) -- ++(0,-0.2) -| ([xshift=6mm]gr.north);
    \draw[route] ([yshift=1.5mm]gr.east) -- node[elabel, above, pos=0.35] {context\\too weak}
      ([yshift=1.5mm]rq.west);
    \draw[loop] ([xshift=6mm]rq.north) -- node[elabel, right] {max 2} ([xshift=6mm]rt.south);
    \draw[route] (gr) -- node[elabel, right] {enough\\context} (ge);
    \draw[route] ([yshift=-1.5mm]gr.east) -- ++(0.45,0) |- (ho.west);
    \node[elabel, above] at (9.4,-8.95) {no relevant clause};
    \draw[route] (ge) -- (og);
    \draw[loop] (og.west) -- ++(-0.3,0) |- ([yshift=-1.5mm]ge.west);
    \node[elabel, left] at (2.85,-9.8) {regenerate,\\max 1};
    \draw[route] (og) -- node[elabel, right] {verified} (fi);
    \draw[route] (og.east) -| (ho.south);
    \node[elabel, above] at (9.4,-10.5) {not verifiable};
    \draw[route] (fi) -- (end);

    % Legend: what colours, borders, arrows and badges mean
    \begin{scope}[yshift=-14.3cm]
      \node[rounded corners=4pt, draw=black!20, fill=black!3, minimum width=15.4cm,
            minimum height=1.55cm] at (6.45,-0.45) {};
      \newcommand\sw[3]{\node[rounded corners=2pt, cat #1, minimum width=0.6cm,
        minimum height=0.32cm] at (#2,0) {}; \node[anchor=west] at (#2+0.4,0) {#3};}
      \sw{model}{-0.6}{language model}
      \sw{sec}{3.3}{check}
      \sw{data}{5.6}{data and search}
      \sw{infra}{9.0}{output}
      \node[rounded corners=2pt, cat infra, dashed, minimum width=0.6cm, minimum height=0.32cm]
        at (11.3,0) {};
      \node[anchor=west] at (11.7,0) {ends the turn};
      \draw[route] (-0.9,-0.85) -- (-0.1,-0.85);
      \node[anchor=west] at (0,-0.85) {route chosen by code};
      \draw[loop] (3.0,-0.85) -- (3.8,-0.85);
      \node[anchor=west] at (3.9,-0.85) {loop with a maximum};
      \fill[black!60] (7.6,-0.85) circle (1.6pt);
      \node[anchor=west] at (7.8,-0.85) {nodes in parallel};
      \node[badge] at (10.9,-0.85) {\includegraphics[height=0.3cm, width=0.34cm,
        keepaspectratio]{logos/openai}};
      \node[anchor=west] at (11.25,-0.85) {technology used};
    \end{scope}
  \end{tikzpicture}
  \caption{The graph of a conversation turn.}
  \label{fig:graph}
\end{figure}


\subsection{Routing Rules}
\label{sec:routing}\index{confidence}\index{policy selection}\index{rewording}

Policy data and clauses are read in parallel, and \code{grade} starts once, after both. Three
decisions are plain rules (Figure~\ref{fig:routing-rules}).

% Figure 2.5: the three decisions of the graph that are rules written in code
\begin{figure}[H]
  \centering
  \begin{tcbraster}[raster columns=3, cardgrid, raster before skip=0pt]
    \begin{tcolorbox}[card]
      \cardtitle{\faFileContract}{Which policy?}
      \begin{tikzpicture}[font=\scriptsize, rung/.style={rounded corners=2pt, draw=accent!60,
          fill=accent!6, text width=3.75cm, align=left, inner sep=2.5pt},
          go/.style={-{Stealth[length=1.6mm]}, accent}]
        \node[rung] (a) at (0,0) {\textbf{1}\enspace the one named (number, plate)};
        \node[rung, below=0.22cm of a] (b) {\textbf{2}\enspace the only active one};
        \node[rung, below=0.22cm of b] (c) {\textbf{3}\enspace the one of the last turn};
        \node[rung, below=0.22cm of c, draw=bad!60, fill=bad!6] (d) {\textbf{?}\enspace ask the
          customer};
        \draw[go] (a) -- (b); \draw[go] (b) -- (c); \draw[go] (c) -- (d);
      \end{tikzpicture}
    \end{tcolorbox}
    \begin{tcolorbox}[card]
      \cardtitle{\faFilter}{Enough context?}
      \begin{tikzpicture}[font=\scriptsize, x=4.3cm]
        \fill[bad!25] (0,0) rectangle (0.55,0.35);
        \fill[good!25] (0.55,0) rectangle (0.80,0.35);
        \fill[good!60] (0.80,0) rectangle (1,0.35);
        \node at (0.275,0.175) {dropped};
        \node at (0.675,0.175) {kept};
        \node[text=white] at (0.9,0.175) {strong};
        \foreach \v in {0, 0.55, 0.80, 1}
          \draw[black!50] (\v,0) -- (\v,-0.08) node[below, font=\tiny] {\v};
        \node[anchor=north, text=black!60] at (0.5,-0.33) {rerank score of a clause};
      \end{tikzpicture}\par
      Confidence shown with the answer:\par
      {\color{good}\faCheck}\ \textbf{high}: one strong clause\par
      {\color{good}\faCheck}\ \textbf{medium}: two kept clauses\par
      {\color{bad}\faTimes}\ otherwise: reword, then operator\par
      {\color{black!55}Thresholds calibrated on the golden set.}
    \end{tcolorbox}
    \begin{tcolorbox}[card]
      \cardtitle{\faRedo}{How many attempts?}
      {\color{accent}\Large\bfseries ×2}\enspace rewordings, with synonyms:\par
      \textit{furto/sottrazione}, \textit{franchigia/scoperto}, \textit{sinistro/evento}\par
      {\color{black!55}(theft, deductible, claim)}\par
      {\color{accent}\Large\bfseries ×1}\enspace regeneration\par
      {\color{bad}\faHeadset}\ then an operator
    \end{tcolorbox}
  \end{tcbraster}
  \caption{The three decisions that are rules, not model outputs. Counters restart at every
    turn.}
  \label{fig:routing-rules}
\end{figure}


\needspace{7cm}
\subsection{State and Context}
\label{sec:state-context}\index{state}\index{context}

The identity never reaches a prompt (Figure~\ref{fig:state-context}).

% Figure 2.6: state (changed by the nodes, saved) versus context (the identity, never in prompts)
\begin{figure}[H]
  \centering
  \begin{tikzpicture}[
      font=\scriptsize,
      panel/.style={rounded corners=4pt, inner sep=5pt, align=left, anchor=north west},
      dest/.style={rounded corners=3pt, align=center, text width=3.6cm, anchor=north},
      ok/.style={-{Stealth[length=2.2mm]}, good, line width=0.9pt},
      no/.style={-{Stealth[length=2.2mm]}, bad, line width=0.9pt, dashed},
    ]
    \node[panel, cat data, text width=7.3cm] (state) at (0,0) {%
      {\small\bfseries \logo{postgresql}\enspace State}\hfill{\color{black!60}saved at every step,
      \acs{aes}}\\[4pt]
      \begin{tabular}{@{}>{\raggedright\arraybackslash}p{2.9cm}@{\hspace{8pt}}
                         >{\raggedright\arraybackslash}p{4.1cm}@{}}
        \textbf{kept between turns} & \textbf{reset at every turn} \\[2pt]
        {\ttfamily messages}\newline {\ttfamily policy\_id}\newline {\ttfamily policy\_meta}
        & {\ttfamily intent, query, policy\_hint, candidates, policy, chunks, context\_ok,
           attempts, gen\_attempts, draft, verdict, handoff\_reason} \\
      \end{tabular}};
    \node[panel, cat sec, text width=6.9cm] (ctx) at (8.45,0) {%
      {\small\bfseries \logo{keycloak}\enspace Context}\hfill{\color{black!60}from the verified
      \acs{jwt}}\\[4pt]
      \begin{tabular}{@{}l@{\hspace{6pt}}>{\raggedright\arraybackslash}p{4.4cm}@{}}
        {\ttfamily customer\_id} & token subject; none for a guest \\
        {\ttfamily access\_token} & {\ttfamily polizze-mcp}, read only; customers only \\
        {\ttfamily channel} & {\ttfamily web} \\
      \end{tabular}};
    \path (state.south) ++(0,-0.9) node[dest, draw=catmodel, fill=catmodel!8] (prompt)
      {\faRobot\enspace\textbf{Prompts}};
    \path (ctx.south) ++(0,-0.9) node[dest, draw=catdata, fill=catdata!10] (mcp)
      {\logo{mcp}\enspace\textbf{Policy server}};
    \draw[ok] (state.south) -- node[right, text=good] {\faCheck} (prompt.north);
    \draw[ok] (ctx.south) -- node[right, text=good] {\faCheck\ token} (mcp.north);
    \draw[no] ([xshift=6mm]ctx.south west) -- node[pos=0.5, fill=white, inner sep=1pt, text=bad]
      {\faTimes\ never} (prompt.north east);
  \end{tikzpicture}
  \caption{The identity lives in the context, which never reaches a prompt: a message like
    \enquote{show me customer 1234's policy} cannot change it. \texttt{input\_guard} resets the
    fields of the turn.}
  \label{fig:state-context}
\end{figure}


\subsection{API and Streaming}
\label{sec:api}\index{streaming}\index{FastAPI}

Two endpoints: the customer's messages and Chatwoot's webhook (Figure~\ref{fig:sequence}). The
client receives progress states, never an unverified draft.

% Figure 2.6: a conversation turn over the API, with progress events and a handover to an operator
\begin{figure}[H]
  \centering
  \begin{tikzpicture}[
      font=\scriptsize,
      actor/.style={rounded corners=3pt, draw=black!40, fill=black!4, align=center,
                    minimum width=2.6cm, minimum height=0.9cm},
      call/.style={-{Stealth[length=2mm]}, black!75, line width=0.6pt},
      event/.style={-{Stealth[length=2mm]}, accent, line width=0.6pt, dashed},
      msg/.style={fill=white, inner sep=1.5pt, align=center},
    ]
    \def\xc{1.3} \def\xa{6.9} \def\xg{10.6} \def\xo{14.4}
    \node[actor] at (\xc,0) {{\color{accent}\faComments}\\Web chat};
    \node[actor] at (\xa,0) {\logo{fastapi}\\Orchestrator \acs{api}};
    \node[actor] at (\xg,0) {\logo{langgraph}\\Graph};
    \node[actor] at (\xo,0) {\logo{chatwoot}\\Chatwoot (operator)};
    \foreach \x in {\xc, \xa, \xg, \xo} \draw[black!30, dashed] (\x,-0.45) -- (\x,-8.7);

    % A normal turn
    \draw[call] (\xc,-1.0) -- node[msg, above] {{\tiny\ttfamily POST /v1/conversazioni/\{conv\_id\}/messaggi}
      \ + \acs{jwt} after login} (\xa,-1.0);
    \draw[call] (\xa,-2.35) -- node[msg, above] {run, with context} (\xg,-2.35);
    \draw[event] (\xg,-2.9) -- (\xc,-2.9);
    \node[msg, anchor=south west, text=accent, yshift=1.5pt] at (\xc+0.15,-2.9) {\code{stato}: progress};
    \draw[event] (\xg,-3.55) -- (\xc,-3.55);
    \node[msg, anchor=south west, text=accent, yshift=1.5pt] at (\xc+0.15,-3.55) {\code{risposta}: after \code{finalize}};

    % A handover
    \draw[black!35, dash dot] (0,-4.05) -- (15.7,-4.05);
    \node[fill=white, text=black!60, font=\scriptsize\itshape] at (8,-4.05) {handover};
    \draw[call] (\xg,-4.6) -- node[msg, above] {ticket, masked summary; pause} (\xo,-4.6);
    \draw[event] (\xg,-5.25) -- (\xc,-5.25);
    \node[msg, anchor=south west, text=accent, yshift=1.5pt] at (\xc+0.15,-5.25) {\code{operatore}: ticket};
    \draw[call] (\xc,-5.9) -- node[msg, above] {message while paused} (\xa,-5.9);
    \draw[call] (\xa,-6.35) -- node[msg, pos=0.75, above] {forwarded} (\xo,-6.35);
    \draw[call] (\xo,-7.1) -- node[msg, above, pos=0.45]
      {{\tiny\ttfamily POST /webhook/chatwoot/conversazione-risolta}, \acs{hmac} verified} (\xa,-7.1);
    \draw[call] (\xa,-7.7) -- node[msg, above] {resume} (\xg,-7.7);
    \draw[event] (\xg,-8.35) -- (\xc,-8.35);
    \node[msg, anchor=south west, text=accent, yshift=1.5pt] at (\xc+0.15,-8.35) {closing note};

    % Legend
    \draw[call] (0,-9.15) -- (0.6,-9.15) node[right, text=black] {call};
    \draw[event] (2.2,-9.15) -- (2.8,-9.15) node[right, text=black] {\acs{sse} event};
  \end{tikzpicture}
  \caption{A turn over the \acs{api}. The thread is keyed by the token subject (or a guest
    session) and the conversation, so no one can resume someone else's.}
  \label{fig:sequence}
\end{figure}


\needspace{7cm}
\subsection{Handover to the Operator}
\label{sec:handover}\index{handover}\index{Chatwoot}

Nine situations end with a person or a fixed reply (Figure~\ref{fig:handover}).

% Figure 2.8: the eight situations in which the assistant does not answer on its own
\begin{figure}[H]
  \centering
  \newcommand\case[2]{\textbf{#1}\newline{\color{black!65}#2}\par\smallskip}
  \begin{tcbraster}[raster columns=3, cardgrid, raster before skip=0pt]
    \begin{tcolorbox}[card, colframe=bad!60]
      \cardtitle{\logo{chatwoot}}{Operator}
      \case{Asks for a person}{\code{classify} $\to$ \code{handoff}}
      \case{Complaint, high priority}{\code{classify} $\to$ \code{handoff}}
      \case{No clause after 2 rewordings}{\code{grade} $\to$ \code{handoff}}
      \case{Unverifiable after 1 regeneration}{\code{output_guard} $\to$ \code{handoff}}
      \case{Writes during a handover}{to Chatwoot; graph stays paused}
    \end{tcolorbox}
    \begin{tcolorbox}[card]
      \cardtitle{\faCommentDots}{Fixed reply or question}
      \case{Blocked message}{\code{input_guard} $\to$ \code{direct_reply}}
      \case{Several policies, none named}{\code{resolve_policy} $\to$ \code{ask_clarification}}
      \case{Personal question, no login}{\code{classify} $\to$ \code{direct_reply}: log in, or
        ask an operator}
    \end{tcolorbox}
    \begin{tcolorbox}[card, colframe=good!60]
      \cardtitle{\faIcon{file-alt}}{Answer + operator offered}
      \case{Coverage of an event}{\code{finalize}: clauses and the notice on claims}
    \end{tcolorbox}
  \end{tcbraster}
  \caption{When the assistant does not answer on its own. Summaries for the operator are masked;
    a rejected draft goes to the operator, never to the customer.}
  \label{fig:handover}
\end{figure}


\subsection{Operator Corrections}
\label{sec:corrections}\index{operator corrections}\index{unanswered questions}

What the operators correct, and what the assistant could not answer, comes back as better
answers, through a human review (Figure~\ref{fig:feedback}).

% Figure 2.9: how the operators' corrections improve the assistant
\begin{figure}[H]
  \centering
  \begin{tikzpicture}[font=\scriptsize,
      fbstep/.style={rounded corners=3pt, text width=2.9cm, minimum height=1.45cm, align=left,
                   inner sep=4pt},
      fbout/.style={rounded corners=3pt, draw=good!60, fill=good!7, text width=3.0cm,
                  minimum height=0.95cm, align=left, inner sep=3pt},
      flow/.style={-{Stealth[length=2.2mm]}, accent, line width=0.8pt}]
    % \fb{name}{x}{category}{icon}{title}{text}
    \newcommand\fb[6]{\node[fbstep, cat #3] (#1) at (#2,0)
      {{\footnotesize\bfseries #4\enspace #5}\\[2pt]#6};}
    \fb{a}{1.6}{infra}{{\color{accent}\faComments}}{Answer}
      {with its sources, the intent and the confidence}
    \fb{o}{5.55}{infra}{\logo{chatwoot}}{Operator}
      {in Chatwoot, marks a wrong answer and writes the right one, with the reason}
    \fb{l}{9.5}{data}{\logo{postgresql}}{Log}
      {corrections and questions left without an answer, masked}
    \fb{r}{13.45}{sec}{{\color{catsec}\faUsers}}{Review}
      {every week the team checks each item: nothing enters unchecked}
    \draw[flow] (a) -- (o);
    \draw[flow] (o) -- (l);
    \draw[flow] (l) -- (r);
    % Approved changes
    \node[fbout] (u1) at (13.45,-2.25) {\faQuestionCircle\ \textbf{FAQ entries} $\to$ the index
      (Figure~\ref{fig:ingestion})};
    \node[fbout] (u2) at (9.5,-2.25) {\faIcon{exchange-alt}\ \textbf{Synonyms} $\to$ the glossary of
      \code{rewrite_query}};
    \node[fbout] (u3) at (5.55,-2.25) {\faDatabase\ \textbf{Test cases} $\to$ the golden set};
    \node[fbout] (u4) at (1.6,-2.25) {\faEdit\ \textbf{Prompt fixes}, tested in \acs{ci} first};
    \begin{scope}[on background layer]
      \node[fit=(u1)(u2)(u3)(u4), rounded corners=5pt, draw=good!50, fill=good!3,
            inner sep=5pt] (box) {};
    \end{scope}
    \draw[flow] (r) -- node[right, text=black!70] {approved} (u1);
    \draw[flow] (box.west) -- ++(-0.15,0) |- node[pos=0.25, right, text=black!70, align=left]
      {better\\answers} (a.west);
  \end{tikzpicture}
  \caption{The operators' corrections, and the questions left without an answer, become better
    answers after a human review.}
  \label{fig:feedback}
\end{figure}


\needspace{7cm}
\subsection{Failure Handling}
\label{sec:failures}\index{failure handling}\index{fallback chain}

Every component has a planned reaction when it fails (Figure~\ref{fig:failures}).

% Figure 2.2: what happens when a component fails
\begin{figure}[H]
  \centering
  \small
  \renewcommand\arraystretch{1.45}
  % \fails{logo}{component}{failure}{reaction}
  \newcommand\fails[4]{\logo[1.1em]{#1} & \textbf{#2} & #3 & {\color{accent}\faArrowRight}
    & #4 \\}
  \begin{tabular}{@{}c@{\hspace{6pt}}l@{\hspace{10pt}}l@{\hspace{6pt}}c@{\hspace{6pt}}l@{}}
    \fails{qwen}{Input guardrail}{no reply in 1\,s}{the turn goes on, the trace records it}
    \fails{openai}{Intent model}{\acs{json} outside the schema}{one more try, then a service message}
    \fails{groq}{Groq quota}{exhausted (429)}{next in the chain: 120b $\to$ 20b $\to$ Qwen3-4B}
    \fails{opensearch}{Search index}{error}{3 tries with growing waits}
    \fails{mcp}{Policy server}{error}{3 tries, then an operator is offered}
    \fails{openai}{Generation}{draft rejected}{one regeneration, then an operator}
    \fails{openai}{Output check}{no reply}{nothing leaves unverified}
    \fails{postgresql}{Audit}{not written}{the answer is not sent}
    \fails{chatwoot}{Handover}{node re-run on resume}{same ticket, never duplicated}
  \end{tabular}
  \caption{What happens when a component fails. Missing data is never estimated by a model.}
  \label{fig:failures}
\end{figure}


\section{Deployment}
\label{sec:deployment}\index{deployment}\index{Docker Compose}\index{latency budget}\index{Oracle Cloud}

Everything runs on one free server, with the two large models on Groq's free \acs{api}
(Figure~\ref{fig:server}). The free quota sets how many questions fit in a day
(Figure~\ref{fig:quota}), and each step of a turn has a time budget
(Figure~\ref{fig:latency}). Memory and times are estimates, to be measured.

% Figure 2.8: the free server (estimated memory of each service) and the free external APIs
\begin{figure}[H]
  \centering
  \begin{tikzpicture}[font=\scriptsize]
    % \svc{row}{logo or icon}{service}{GB}{style}
    \newcommand\svc[5]{%
      \node at (0.3,-#1*0.55) {#2};
      \node[anchor=west] at (0.65,-#1*0.55) {#3};
      \fill[#5] (5.2,-#1*0.55-0.16) rectangle ({5.2+#4*0.55},-#1*0.55+0.16);
      \node[anchor=west, text=black!70] at ({5.25+#4*0.55},-#1*0.55) {#4\,GB};}
    \newcommand\lgx[1]{\includegraphics[height=0.32cm, width=0.4cm, keepaspectratio]{logos/#1}}
    \tikzset{on/.style={fill=catdata!70}, demand/.style={pattern=north east lines,
      pattern color=catdata!80}}
    % Server frame
    \draw[rounded corners=4pt, draw=catinfra, line width=0.7pt, fill=catinfra!4]
      (-0.1,0.95) rectangle (9.4,-6.25);
    \node[anchor=west, font=\footnotesize\bfseries] at (0.05,0.55)
      {\lgx{oracle}\enspace Free Oracle Cloud server \textmd{\color{black!60} 4 Arm cores,
      24\,GB}};
    \svc{0.4}{\lgx{opensearch}}{OpenSearch (1\,GB heap)}{2.0}{on}
    \svc{1.4}{\lgx{keycloak}}{Keycloak}{0.8}{on}
    \svc{2.4}{\lgx{chatwoot}}{Chatwoot + Valkey}{2.2}{on}
    \svc{3.4}{\lgx{postgresql}}{PostgreSQL}{0.5}{on}
    \svc{4.4}{\lgx{langgraph}}{API, graph, policy server}{0.6}{on}
    \svc{5.4}{\lgx{qwen}}{Three Qwen3 0.6B (int8)}{2.0}{on}
    \svc{6.4}{{\color{catinfra}\faServer}}{Other services on it}{0.5}{on}
    \svc{7.4}{\lgx{ollama}}{Qwen3-4B fallback}{3.0}{demand}
    \svc{8.4}{\lgx{docling}}{Docling, when ingesting}{3.5}{demand}
    \draw[black!30] (0.1,-5.15) -- (9.2,-5.15);
    \node[anchor=west, font=\footnotesize] at (0.1,-5.55)
      {\textbf{Always on} $\approx$ 8.6\,GB \quad \textbf{Peak} $\approx$ 15\,GB \quad
       \textbf{Free} $\geq$ 9\,GB};
    \node[anchor=west, text=black!55] at (0.1,-5.95) {Estimates, to be measured on the server.};

    % External free APIs
    \node[anchor=north west, rounded corners=4pt, draw=catmodel, fill=catmodel!6, text width=5.6cm,
          inner sep=5pt, align=left] at (9.9,0.95) {%
      {\footnotesize\bfseries \lgx{groq}\enspace Groq, free plan}\\[3pt]
      \lgx{openai}\ gpt-oss-120b \textbullet\ gpt-oss-20b\\[2pt]
      each: 30 requests/min, 1{,}000/day,\\ 8K tokens/min, 200K tokens/day};
    \node[anchor=north west, rounded corners=4pt, draw=catmodel, fill=catmodel!6, text width=5.6cm,
          inner sep=5pt, align=left] at (9.9,-1.75) {%
      {\footnotesize\bfseries \lgx{gemini}\enspace Google AI Studio, free}\\[3pt]
      Gemini Flash, only as the judge of the tests};
    \node[anchor=north west, text width=5.6cm, align=left, text=black!60] at (9.9,-3.3)
      {Only what the answer needs is sent, with the identifiers masked.};
  \end{tikzpicture}
  \caption{Where everything runs: one free server, plus two free \acsp{api}~\cite{groq2026limits}.
    Hatched: memory used only when needed.}
  \label{fig:server}
\end{figure}

% Figure 2.9: what one question costs in free quota, and how many questions fit in a day
\begin{figure}[H]
  \centering
  \begin{tikzpicture}[font=\scriptsize]
    \newcommand\lgx[1]{\includegraphics[height=0.32cm, width=0.4cm, keepaspectratio]{logos/#1}}
    % One question: calls and tokens (estimated)
    \node[anchor=west, font=\footnotesize\bfseries] at (0,0.6) {One question};
    % \call{row}{logo}{step}{tokens in thousands}{style}{where}
    \newcommand\call[6]{%
      \node at (0.25,-#1*0.55) {\lgx{#2}};
      \node[anchor=west, font=\scriptsize\ttfamily] at (0.55,-#1*0.55) {#3};
      \fill[#5] (3.4,-#1*0.55-0.15) rectangle ({3.4+#4*0.6},-#1*0.55+0.15);
      \node[anchor=west, text=black!70] at ({3.45+#4*0.6},-#1*0.55) {#6};}
    \call{0}{qwen}{input\_guard}{0}{white}{on the server: no quota}
    \call{1}{openai}{classify}{1.5}{catmodel!35}{$\approx$1.5K tokens, 20b}
    \call{2}{openai}{rewrite\_query}{0.5}{catmodel!35}{0--2 $\times$ 0.5K, 20b}
    \call{3}{openai}{generate}{4}{catmodel}{$\approx$4K tokens, 120b}
    \call{4}{openai}{output\_guard}{4}{catmodel!35}{$\approx$4K tokens, 20b}
    \call{5}{qwen}{embed, rerank}{0}{white}{on the server: no quota}

    % Daily capacity
    \begin{scope}[xshift=10.2cm]
      \node[anchor=west, font=\footnotesize\bfseries] at (0,0.6) {Questions a day};
      \node[anchor=west] at (0,-0.1) {\lgx{openai}\ gpt-oss-120b: 200K / 4K};
      \fill[catmodel] (0,-0.75) rectangle (5.0,-0.45);
      \node[anchor=east, text=white, font=\scriptsize\bfseries] at (5.0,-0.6) {$\approx$50};
      \node[anchor=west] at (0,-1.25) {\lgx{openai}\ gpt-oss-20b: 200K / 5.5K};
      \fill[catmodel!35] (0,-1.9) rectangle (3.6,-1.6);
      \node[anchor=east, font=\scriptsize\bfseries] at (3.6,-1.75) {$\approx$36};
      \node[anchor=west, text width=5.4cm, align=left] at (0,-2.65)
        {\textbf{About 35 questions a day}, set by the 20b. Then the chain falls back to Qwen3-4B
         on the server, slower but without limits.};
    \end{scope}
  \end{tikzpicture}
  \caption{The free quota of one question (estimated tokens) and the questions it allows each day.
    During development, a cache keeps repeated calls from using quota.}
  \label{fig:quota}
\end{figure}

% Figure 2.10: time budget of every step of a turn (1.35 cm per second), as design targets
\begin{figure}[H]
  \centering
  \begin{tikzpicture}[font=\scriptsize, x=1.35cm, y=0.95cm]
    \def\xo{3.45}  % the time axis starts here, in units of 1.35 cm
    % \step{row}{logo}{node}{note}{from s}{to s}{category}
    \newcommand\step[7]{%
      \node at (0.13,-#1) {\includegraphics[height=0.4cm, width=0.45cm, keepaspectratio]{logos/#2}};
      \node[anchor=west, align=left, text width=3.8cm] at (0.3,-#1)
        {{\ttfamily #3}\\[-1pt]{\tiny\color{black!60}#4}};
      \draw[cat #7, rounded corners=1.5pt] (\xo+#5,-#1-0.3) rectangle (\xo+#6,-#1+0.3);}
    \step{1}{qwen}{input\_guard}{0.6B model on the CPU}{0}{0.5}{sec}
    \step{2}{openai}{classify}{short \acs{json} on Groq}{0.5}{1.0}{model}
    \step{3}{mcp}{resolve\_policy}{core insurance system \acsp{api}}{1.0}{1.2}{data}
    \step{4}{mcp}{fetch\_policy\_data}{in parallel with retrieve}{1.2}{1.4}{data}
    \step{5}{opensearch}{retrieve}{rerank of 40 on the CPU: the slower}{1.2}{3.2}{data}
    \step{6}{openai}{generate}{about 300 tokens on Groq}{3.2}{4.7}{model}
    \step{7}{openai}{output\_guard}{validator + gpt-oss-20b}{4.7}{5.7}{sec}
    % Durations
    \foreach \r/\a/\b/\t in {1/0/0.5/0.5, 2/0.5/1.0/0.5, 3/1.0/1.2/0.2, 4/1.2/1.4/0.2,
                             5/1.2/3.2/2.0, 6/3.2/4.7/1.5, 7/4.7/5.7/1.0}
      \node[anchor=west, font=\tiny] at (\xo+\b,-\r) {\t\,s};
    % Axis
    \draw[black!50] (\xo,-7.6) -- (\xo+8.2,-7.6);
    \foreach \s in {0,...,8}
      \draw[black!50] (\xo+\s,-7.6) -- (\xo+\s,-7.75) node[below] {\s\,s};
    % Targets
    \draw[bad, dashed, line width=0.8pt] (\xo+8,0.2) -- (\xo+8,-7.6);
    \node[anchor=south east, text=bad, font=\scriptsize\bfseries] at (\xo+8.05,0.2)
      {target: under 8\,s};
    \draw[good, dashed, line width=0.8pt] (\xo+1,0.2) -- (\xo+1,-7.6);
    \node[anchor=south, text=good, font=\scriptsize\bfseries, align=center] at (\xo+1,0.2)
      {first progress state within 1\,s};
    \draw[decorate, decoration={brace, mirror, amplitude=4pt}, black!60]
      (\xo,-8.45) -- node[below=4pt] {sum of the budgets: 5.7\,s} (\xo+5.7,-8.45);
  \end{tikzpicture}
  \caption{Time budget per step, as design targets to be measured (p95: 95\% of the requests
    within the time).}
  \label{fig:latency}
\end{figure}

